\chapter{Hỗ trợ --- Yêu cầu thông tin --- Kiến thức --- Thông báo}

Chương này gồm \textbf{22 đường dẫn} gom thành \textbf{19 mục}, thuộc bốn mô-đun:
\rt{wujia_portal_support} (5), \rt{wujia_portal_info_request} (5),
\rt{wujia_portal_knowledge} (4), \rt{wujia_portal_notification} (8).

\textbf{Bốn điểm cần biết trước khi đọc:}

\begin{enumerate}[leftmargin=*]
  \item \textbf{Phiếu Hỗ trợ lọc theo NGƯỜI TẠO, không theo cửa hàng.} Đây là phân hệ duy nhất
        trong portal làm như vậy. Đồng nghiệp cùng cửa hàng \textbf{không thấy} phiếu của nhau;
        Chủ tiệm \textbf{không thấy} phiếu nhân viên mình gửi. Trong khi đó Yêu cầu thông tin
        (ngay chương này) lại lọc theo cửa hàng nên cả cửa hàng cùng thấy. Hai phân hệ giống
        nhau về hình thức nhưng ngược nhau về phạm vi.
  \item \textbf{Kiến thức không phân biệt cửa hàng và không phân biệt vai trò.} Mọi tài khoản
        portal thấy đúng cùng một tập bài viết. Đây là phân hệ duy nhất như vậy.
  \item \textbf{Thông báo ghi ``đã đọc'' theo bộ ba người + thông báo + cửa hàng.} Cùng một
        người, đọc ở cửa hàng A rồi bấm sang cửa hàng B thì thông báo đó \textbf{lại thành chưa
        đọc}. UAT hiện có \uat{6} dấu đã đọc cũ \textbf{chưa gắn cửa hàng} (dữ liệu trước khi đổi
        cách ghi) --- những dấu này không khớp với bất kỳ cửa hàng nào.
  \item \textbf{Yêu cầu cập nhật thông tin chưa từng được dùng.} UAT có \uat{0} bản ghi. Mọi mô
        tả trong phần đó là đọc từ mã nguồn, chưa có số thật để đối chiếu.
\end{enumerate}

Quy ước chung toàn portal xem mục~\ref{sec:quy-uoc-chung}.

% =============================================================================
\section{Màn Danh sách phiếu hỗ trợ}
\rt{/portal/support}
\medskip

\mvao

Mục \emph{Hỗ trợ} trên thanh điều hướng. Tham số: \rt{page}, \rt{state}, \rt{q}.

\mai

\begin{itemize}[leftmargin=*]
  \item Phải đăng nhập. \textbf{Không chặn vai trò.}
  \item Chỉ thấy phiếu \textbf{do chính tài khoản mình tạo} (\rt{created_by_id = tôi}).
        \textbf{Không} lọc theo cửa hàng --- nghĩa là người rời cửa hàng vẫn thấy phiếu cũ của
        mình, còn người mới vào cửa hàng không thấy gì.
  \item Phiếu bị Ngô Gia tắt cờ hiện trên portal thì \textbf{biến mất hoàn toàn}, kể cả với
        người đã tạo ra nó.
  \item Thực tế UAT: \uat{15} phiếu, \uat{0} phiếu do tài khoản cửa hàng thật tạo --- toàn bộ
        do tài khoản nội bộ tạo. Vì vậy \textbf{mọi tài khoản cửa hàng trên UAT đều mở ra màn
        rỗng}, và bản chụp màn hình phải làm bằng tài khoản Administrator.
\end{itemize}

\mhien

\begin{hienbang}
Danh sách phiếu &
  \rt{wujia.support.ticket} &
  Phiếu do tôi tạo, còn bật cờ hiện trên portal, và \textbf{không phải phiếu đã huỷ}
  (\rt{state != 'cancelled'}) &
  Ngày tạo mới nhất trước; \textbf{20} phiếu mỗi trang, \textbf{cố định} &
  \uat{15} phiếu toàn hệ, \uat{2} đã huỷ (bị giấu) \\ \addlinespace
Huy hiệu trạng thái &
  \rt{state} &
  Sáu nhãn: \emph{Mới} · \emph{Đang xử lý} · \emph{Chờ phản hồi} · \emph{Đã giải quyết} ·
  \emph{Đã đóng} · \emph{Đã huỷ} (nhãn cuối chỉ gặp ở màn chi tiết) &
  --- &
  \uat{4} Mới · \uat{3} Đang xử lý · \uat{1} Chờ phản hồi · \uat{3} Đã giải quyết ·
  \uat{2} Đã đóng \\ \addlinespace
Huy hiệu mức độ &
  \rt{priority} &
  Hai mức: \emph{Bình thường} và \emph{Khẩn} &
  --- &
  \uat{4}/15 phiếu đánh dấu Khẩn \\ \addlinespace
Cột danh mục &
  \rt{wujia.support.category} &
  Danh mục do Ngô Gia cấu hình &
  --- &
  \uat{7} danh mục đang bật \\
\end{hienbang}

\mloc

\begin{itemize}[leftmargin=*]
  \item \textbf{Trạng thái} --- một trong sáu nhãn; giá trị lạ bị bỏ qua.
  \item \textbf{Từ khoá} --- khớp một phần theo \textbf{mã phiếu} hoặc \textbf{tiêu đề}. Không
        tìm trong nội dung mô tả, không tìm trong các lời nhắn.
  \item \textbf{Không có bộ lọc ngày} và \textbf{không đổi được số dòng mỗi trang} --- khác mọi
        màn danh sách khác trong portal.
\end{itemize}

\mkiem{Tạo phiếu}

Không kiểm gì ở màn danh sách.

\mghi

\textbf{Không ghi gì.}

\mhoi

\begin{ask}
\begin{enumerate}[leftmargin=*]
  \item \textbf{Phiếu chỉ người tạo mới thấy.} Nhân viên nghỉ việc thì phiếu đi theo tài khoản
        đó; Chủ tiệm không nắm được cửa hàng mình đã gửi những gì. Anh muốn đổi sang
        \textbf{lọc theo cửa hàng} như các phân hệ còn lại không? Đây là điểm đã nêu ở bản tài
        liệu cũ và tới nay chưa chốt.
  \item \textbf{Không có bộ lọc ngày và cỡ trang cố định 20} --- lệch với năm màn danh sách
        khác. Có cần đồng bộ không?
  \item \textbf{Phiếu Đã huỷ bị giấu khỏi danh sách nhưng màn chi tiết vẫn mở được} nếu còn giữ
        đường dẫn. UAT có \uat{2} phiếu như vậy.
\end{enumerate}
\end{ask}

% =============================================================================
\section{Màn Tạo phiếu hỗ trợ}
\rt{/portal/support/new}
\medskip

\mvao

Nút \emph{Tạo phiếu} ở màn danh sách. Cùng một đường dẫn dùng cho cả mở form lẫn gửi form.

\mai

Phải đăng nhập và \textbf{truy cập được ít nhất một cửa hàng}; không thì bị đưa thẳng về danh
sách (không có thông báo gì). Không chặn vai trò.

\mhien

\begin{hienbang}
Ô \emph{Cửa hàng} &
  \rt{wujia.franchise.management} &
  Mọi cửa hàng tôi truy cập được &
  --- & --- \\ \addlinespace
Ô \emph{Danh mục} &
  \rt{wujia.support.category} &
  Danh mục đang bật (\rt{active = true}) &
  Theo số thứ tự, rồi tên &
  \uat{7} danh mục \\ \addlinespace
Ô \emph{Mức độ} &
  --- &
  Hai lựa chọn: \emph{Bình thường} (mặc định) và \emph{Khẩn} &
  --- & --- \\ \addlinespace
Khối tệp đính kèm &
  \rt{ir.attachment} &
  Tối đa \textbf{6 tệp}, mỗi tệp $\leq$ \textbf{5 MB}; chỉ nhận PNG, JPEG, WEBP, PDF &
  --- & --- \\
\end{hienbang}

\mloc

\begin{itemize}[leftmargin=*]
  \item \textbf{Tiêu đề} --- bắt buộc. Hệ thống nhận cả tên ô cũ (\rt{subject}) lẫn tên ô mới
        (\rt{title}) để không làm hỏng liên kết cũ.
  \item \textbf{Mô tả} --- không bắt buộc, không giới hạn độ dài ở tầng portal.
  \item \textbf{Mức độ} --- giá trị lạ tự hạ về \emph{Bình thường}.
\end{itemize}

\mkiem{Gửi phiếu}

\begin{kiembang}
1 & Cửa hàng và danh mục gửi lên là số & Quay lại form kèm mã lỗi \rt{invalid_input} \\
2 & Có tiêu đề, có cửa hàng, có danh mục & Quay lại form kèm \rt{missing_fields} \\
3 & Cửa hàng \textbf{nằm trong danh sách của tôi} & Quay lại form kèm \rt{invalid_franchise} \\
4 & Danh mục còn đang bật & Quay lại form kèm \rt{invalid_input} \\
5 & Tệp đính kèm đúng số lượng, dung lượng, định dạng &
    Phiếu vừa tạo \textbf{bị xoá hẳn}, quay lại form kèm \rt{invalid_attachment} \\
\end{kiembang}

Bốn bước đầu chạy \textbf{trước} khi ghi; bước 5 chạy \textbf{sau} khi đã ghi phiếu, nên hệ
thống phải xoá lại phiếu. Các mã lỗi trên được trả về dưới dạng tham số trên thanh địa chỉ và
màn dịch ra câu tiếng Việt.

\mghi

\begin{itemize}[leftmargin=*]
  \item Tạo một phiếu \rt{wujia.support.ticket}, mã tự sinh (UAT đang ở \emph{WJ-TK/26/00016}),
        trạng thái \textbf{Mới}, cờ hiện trên portal \textbf{bật sẵn}.
  \item Người tạo = người đang đăng nhập; cửa hàng = cửa hàng chọn trong form.
  \item Tạo các bản ghi tệp đính kèm gắn vào phiếu.
  \item \textbf{Không} gửi email, \textbf{không} tạo việc cần làm cho ai. Phiếu chỉ hiện trong
        ERP để Ngô Gia tự vào xem.
\end{itemize}

\mhoi

\begin{ask}
\begin{enumerate}[leftmargin=*]
  \item \textbf{Định dạng tệp chỉ kiểm theo lời khai của trình duyệt}, không đọc nội dung thật
        --- khác với màn Đổi trả (chương 6) đã đọc nội dung. Đây là mục đã nằm trong danh sách
        chuẩn hoá F1 của dev; ghi ở đây để BA biết, không cần ra task riêng.
  \item \textbf{Gửi phiếu xong không ai được báo.} Không email, không thông báo trong ERP. Ngô
        Gia phải chủ động mở danh sách. Anh có muốn gắn thông báo cho người phụ trách không?
  \item \textbf{Chưa có giới hạn số phiếu gửi trong ngày} --- một tài khoản bấm liên tục sẽ tạo
        liên tục.
\end{enumerate}
\end{ask}

% =============================================================================
\section{Màn Chi tiết phiếu hỗ trợ}
\rt{/portal/support/<int:ticket_id>}
\medskip

\mvao

Bấm một dòng ở màn danh sách, hoặc được đưa tới ngay sau khi tạo phiếu.

\mai

Phiếu phải \textbf{do chính tôi tạo} và \textbf{còn bật cờ hiện trên portal}. Sai thì bị đưa về
danh sách, không có thông báo. \textbf{Lưu ý}: màn này \textbf{không} loại phiếu đã huỷ --- phiếu
huỷ bị giấu ở danh sách nhưng vẫn mở được bằng đường dẫn.

\mhien

\begin{hienbang}
Đầu phiếu &
  \rt{wujia.support.ticket} &
  Mã phiếu, tiêu đề, danh mục, mức độ, trạng thái, cửa hàng, ngày tạo &
  --- & --- \\ \addlinespace
Nội dung &
  \rt{description} &
  Mô tả người dùng gõ lúc tạo &
  --- & --- \\ \addlinespace
Khối trao đổi &
  \rt{mail.message} &
  \textbf{Mọi lời nhắn} trên phiếu (loại bỏ các dòng nhật ký hệ thống). Nghĩa là lời Ngô Gia gõ
  vào ô trao đổi trong ERP \textbf{hiện thẳng} cho cửa hàng. Ô \emph{Ghi chú nội bộ} thì không. &
  Cũ trước, mới sau &
  --- \\ \addlinespace
Tệp đính kèm &
  \rt{ir.attachment} &
  Tệp gắn vào phiếu, tải qua đường dẫn có kiểm quyền &
  --- &
  \uat{1} tệp trên toàn bộ 15 phiếu \\
\end{hienbang}

\mloc

Chỉ một ô: khung soạn lời nhắn.

\mkiem{Gửi lời nhắn}

Xem mục kế tiếp.

\mghi

Mở màn \textbf{không ghi gì}.

\mhoi

\begin{ask}
\begin{enumerate}[leftmargin=*]
  \item \textbf{Ranh giới nội bộ / công khai nằm ở chỗ nhân sự Ngô Gia chọn ô nào.} Gõ vào ô
        trao đổi thì cửa hàng đọc được; muốn ghi riêng phải dùng ô \emph{Ghi chú nội bộ} hoặc
        chức năng ghi chú. Cần dặn lại đội vận hành điểm này.
  \item \textbf{Phiếu Đã đóng / Đã huỷ vẫn nhận được lời nhắn mới} --- không có chốt chặn theo
        trạng thái.
\end{enumerate}
\end{ask}

% =============================================================================
\section{Gửi lời nhắn vào phiếu hỗ trợ}
\rt{/portal/support/<int:ticket_id>/reply}
\medskip

\mvao

Không phải màn --- là hành động của nút \emph{Gửi} trong khung trao đổi ở màn chi tiết.

\mai

Cùng điều kiện với màn chi tiết: phiếu do tôi tạo và còn hiện trên portal.

\mhien

Không hiện gì --- xong việc thì quay lại chính màn chi tiết.

\mloc

Đúng một ô nội dung. \textbf{Không} đính kèm được tệp trong lời nhắn.

\mkiem{Gửi}

\begin{kiembang}
1 & Phiếu tồn tại, do tôi tạo, còn hiện trên portal & Đưa về danh sách \\
2 & Nội dung \textbf{không rỗng} sau khi bỏ khoảng trắng &
    \textbf{Không báo lỗi} --- lặng lẽ quay về màn chi tiết, không ghi gì \\
\end{kiembang}

\mghi

\begin{itemize}[leftmargin=*]
  \item Ghi một lời nhắn vào phiếu, \textbf{đứng tên chính người dùng} (không phải tên hệ
        thống), loại \emph{bình luận} nên Ngô Gia nhận được trong ERP.
  \item Việc ghi lời nhắn \textbf{kích hoạt phép tính thời gian phản hồi} của phiếu --- dùng
        cho báo cáo nội bộ.
  \item \textbf{Không} đổi trạng thái phiếu. Phiếu đang \emph{Chờ phản hồi} mà cửa hàng trả lời
        thì vẫn nằm nguyên ở \emph{Chờ phản hồi} cho tới khi Ngô Gia tự đổi.
\end{itemize}

\mhoi

\begin{ask}
\begin{enumerate}[leftmargin=*]
  \item \textbf{Trả lời không tự đổi trạng thái.} Anh có muốn phiếu \emph{Chờ phản hồi} tự về
        \emph{Đang xử lý} khi cửa hàng trả lời không?
  \item \textbf{Không đính kèm được ảnh trong lời nhắn} --- cửa hàng muốn gửi thêm ảnh chụp
        phải mở phiếu mới.
\end{enumerate}
\end{ask}

% =============================================================================
\section{Tải tệp đính kèm của phiếu hỗ trợ}
\rt{/portal/support/<int:ticket_id>/attachment/<int:att_id>}
\medskip

\mvao

Bấm tên tệp ở màn chi tiết.

\mai

\begin{kiembang}
1 & Phiếu tồn tại, do tôi tạo, còn hiện trên portal & Báo không tìm thấy \\
2 & Tệp \textbf{thuộc đúng phiếu đó} &
    Từ chối truy cập --- không mượn được đường dẫn này để đọc tệp khác trong hệ thống \\
\end{kiembang}

\mhien

Trả về tệp ở chế độ \textbf{tải xuống} (khác Đổi trả --- bên đó mở thẳng trong trình duyệt).

\mloc

Không có.

\mkiem{tên tệp}

Xem bảng trên.

\mghi

\textbf{Không ghi gì.}

\mhoi

\begin{ask}
Không có điểm cần chốt.
\end{ask}

% =============================================================================
\section{Màn Danh sách yêu cầu cập nhật thông tin}
\rt{/portal/info-request}
\medskip

\mvao

Mục \emph{Yêu cầu cập nhật thông tin} trên thanh điều hướng. Tham số: \rt{page},
\rt{state}, \rt{request_type}, \rt{page_size}.

\mai

\begin{itemize}[leftmargin=*]
  \item Phải đăng nhập. \textbf{Màn danh sách không chặn vai trò} --- Nhân viên xem được;
        chỉ màn \emph{tạo} mới chặn.
  \item Phạm vi là \textbf{mọi cửa hàng tôi truy cập được}, không lọc theo người tạo --- ngược
        hẳn với phân hệ Hỗ trợ ở đầu chương.
  \item Không truy cập được cửa hàng nào thì màn rỗng.
\end{itemize}

\mhien

\begin{hienbang}
Danh sách yêu cầu &
  \rt{wujia.info.update.request} &
  Yêu cầu thuộc mọi cửa hàng của tôi &
  Ngày tạo mới nhất trước; \textbf{20} yêu cầu mỗi trang (10/20/50) &
  \uat{0} bản ghi --- phân hệ chưa được dùng lần nào \\ \addlinespace
Huy hiệu trạng thái &
  \rt{state} &
  Năm nhãn: \emph{Nháp} · \emph{Đã gửi} · \emph{Đang xem} · \emph{Đã duyệt} · \emph{Từ chối}.
  \textbf{Không có} nhãn ``Đã huỷ''. &
  --- & --- \\ \addlinespace
Cột loại thông tin &
  \rt{request_type} &
  Bảy loại cố định trong mã nguồn: Địa chỉ · Số điện thoại · Email · Tên chủ cửa hàng ·
  Thông tin ngân hàng · Người đại diện · Khác &
  --- & --- \\
\end{hienbang}

\mloc

Lọc theo \textbf{trạng thái} và \textbf{loại thông tin}; giá trị lạ bị bỏ qua. Số dòng mỗi
trang 10/20/50. \textbf{Không} có ô tìm kiếm, \textbf{không} có bộ lọc ngày.

\mkiem{Tạo yêu cầu}

Xem màn kế tiếp.

\mghi

\textbf{Không ghi gì.}

\mhoi

\begin{ask}
\begin{enumerate}[leftmargin=*]
  \item \textbf{Nhân viên xem được mọi yêu cầu của cửa hàng} (kể cả yêu cầu đổi thông tin ngân
        hàng) nhưng \textbf{không tạo được}. Đúng ý anh chứ?
  \item \textbf{Không có nhãn ``Đã huỷ''} --- yêu cầu người dùng tự huỷ sẽ hiện thành
        \emph{Từ chối}; xem mục Huỷ yêu cầu bên dưới.
  \item \textbf{Bảy loại thông tin đang cố định trong mã nguồn}, muốn thêm loại phải sửa mã.
        Anh xác nhận danh sách này đủ.
\end{enumerate}
\end{ask}

% =============================================================================
\section{Màn Tạo yêu cầu cập nhật thông tin}
\rt{/portal/info-request/new}
\medskip

\mvao

Nút \emph{Tạo yêu cầu} ở màn danh sách. Cùng đường dẫn cho cả mở form lẫn gửi form.

\mai

\begin{itemize}[leftmargin=*]
  \item Phải truy cập được ít nhất một cửa hàng.
  \item \textbf{Chỉ Chủ tiệm / Quản lý}. Nhân viên nhận \textbf{trang lỗi 403 thô của máy chủ},
        \textbf{không} phải trang thông báo thân thiện như phân hệ Công nợ. Hai phân hệ cùng
        chặn vai trò nhưng chặn ra hai kiểu khác nhau.
  \item Vai trò xét theo \textbf{mức cao nhất trên MỌI cửa hàng} của tài khoản --- giống màn Báo
        cáo (chương 5), khác Công nợ (chương 6, xét tại cửa hàng đang chọn).
\end{itemize}

\mhien

\begin{hienbang}
Ô \emph{Cửa hàng} &
  \rt{wujia.franchise.management} &
  Mọi cửa hàng tôi truy cập được &
  --- & --- \\ \addlinespace
Ô \emph{Loại thông tin} &
  --- &
  Bảy loại cố định &
  Theo thứ tự khai trong mã &
  --- \\ \addlinespace
Ô \emph{Giá trị hiện tại} &
  \rt{wujia.franchise.management} &
  Đọc thẳng từ cửa hàng theo loại thông tin đang chọn; gọi ngầm mỗi lần đổi ô
  (xem mục \emph{Lấy giá trị hiện tại}) &
  --- & --- \\ \addlinespace
Khối tệp đính kèm &
  \rt{ir.attachment} &
  Tối đa \textbf{6 tệp}, mỗi tệp $\leq$ \textbf{5 MB}; PNG, JPEG, WEBP, PDF &
  --- & --- \\
\end{hienbang}

\mloc

\begin{itemize}[leftmargin=*]
  \item \textbf{Giá trị mới} --- bắt buộc, cắt còn 5.000 ký tự.
  \item \textbf{Tên trường cần đổi} --- chỉ hiện khi chọn loại \emph{Khác}; bắt buộc khi đó.
        Đây là \textbf{tên kỹ thuật} của trường, không phải tên tiếng Việt.
  \item \textbf{Ghi chú} --- cắt còn 2.000 ký tự.
  \item \textbf{Mức độ} --- Bình thường / Khẩn; giá trị lạ hạ về Bình thường.
  \item \textbf{Hai nút}: \emph{Lưu nháp} và \emph{Gửi}.
\end{itemize}

\mkiem{Gửi yêu cầu}

\begin{kiembang}
1 & Truy cập được ít nhất một cửa hàng & Đưa về danh sách \\
2 & Vai trò Chủ tiệm / Quản lý (tính trên mọi cửa hàng) & Trang lỗi 403 của máy chủ \\
3 & Cửa hàng gửi lên là số và nằm trong danh sách của tôi &
    ``Cửa hàng không hợp lệ.'' / ``Cửa hàng không truy cập được.'' \\
4 & Loại thông tin nằm trong bảy loại & ``Loại thông tin không hợp lệ.'' \\
5 & Có nhập giá trị mới & ``Vui lòng nhập giá trị mới.'' \\
6 & Chọn \emph{Khác} thì phải nhập tên trường & ``Khi chọn 'Khác' phải nhập tên field.'' \\
7 & Ghi được bản ghi &
    ``Không thể gửi yêu cầu. Vui lòng kiểm tra lại thông tin và thử lại.'' \\
8 & Tệp đính kèm hợp lệ & Bản ghi vừa tạo \textbf{bị xoá hẳn}, quay lại form \\
\end{kiembang}

\mghi

\begin{itemize}[leftmargin=*]
  \item Tạo một bản ghi \rt{wujia.info.update.request}, mã tự sinh, trạng thái \textbf{Nháp}.
  \item Bấm \emph{Gửi} thì chuyển sang \textbf{Đã gửi}, ghi mốc thời gian gửi và một dòng vào
        khung trao đổi của bản ghi.
  \item \textbf{Không} sửa gì trên hồ sơ cửa hàng --- đây chỉ là đề nghị.
\end{itemize}

\mhoi

\begin{ask}
\begin{enumerate}[leftmargin=*]
  \item \textbf{Quyền xét trên mọi cửa hàng nhưng gửi được cho từng cửa hàng.} Một người là
        Quản lý cửa hàng A và Nhân viên cửa hàng B vẫn gửi được yêu cầu đổi thông tin
        \textbf{của cửa hàng B}. Anh muốn siết theo đúng cửa hàng như Công nợ không?
  \item \textbf{Nhân viên bị chặn bằng trang lỗi kỹ thuật 403}, không phải trang thông báo.
        Cần làm giống Công nợ.
  \item \textbf{Loại ``Khác'' bắt người dùng gõ tên kỹ thuật của trường} --- không ai ngoài dev
        biết gõ gì. Anh muốn đổi thành ô mô tả tự do không?
  \item \textbf{Duyệt yêu cầu KHÔNG tự cập nhật hồ sơ cửa hàng.} Bấm duyệt chỉ đổi trạng thái;
        nhân sự Ngô Gia vẫn phải tự vào sửa tay. Anh xác nhận giữ vậy, hay muốn duyệt là áp
        luôn?
  \item \textbf{Phiếu nháp ở đây gửi được} (khác Đổi trả) --- nhưng chỉ khi tạo lại từ đầu; màn
        chi tiết cũng không có nút gửi. Cần rà lại cho thống nhất.
\end{enumerate}
\end{ask}

% =============================================================================
\section{Màn Chi tiết yêu cầu cập nhật thông tin}
\rt{/portal/info-request/<int:req_id>}
\medskip

\mvao

Bấm một dòng ở màn danh sách, hoặc sau khi tạo xong.

\mai

Yêu cầu phải thuộc \textbf{một trong các cửa hàng của tôi} --- \textbf{không} đòi phải là người
tạo. Sai thì về danh sách. Không chặn vai trò.

\mhien

\begin{hienbang}
Đầu yêu cầu &
  \rt{wujia.info.update.request} &
  Mã, cửa hàng, loại thông tin, trạng thái, mức độ, ngày tạo, ngày gửi &
  --- & --- \\ \addlinespace
Giá trị cũ &
  \rt{wujia.franchise.management} &
  \textbf{Đọc lại từ cửa hàng ngay lúc mở màn} --- không phải ảnh chụp lúc gửi. Ngô Gia sửa hồ
  sơ xong thì ô ``giá trị cũ'' của yêu cầu cũ \textbf{đổi theo}, và trông như cũ đã bằng mới. &
  --- & --- \\ \addlinespace
Giá trị mới &
  \rt{new_value} &
  Chữ người dùng đã gõ &
  --- & --- \\ \addlinespace
Tệp đính kèm · lý do từ chối &
  \rt{ir.attachment} · \rt{refuse_reason} &
  Tệp minh chứng và lý do khi bị từ chối &
  --- & --- \\
\end{hienbang}

\mloc

Không có ô lọc.

\mkiem{Huỷ yêu cầu}

Xem mục kế tiếp.

\mghi

Mở màn \textbf{không ghi gì}.

\mhoi

\begin{ask}
\textbf{``Giá trị cũ'' không phải ảnh chụp.} Mô tả trong mã nguồn ghi là ảnh chụp lúc mở form,
nhưng thực tế nó đọc lại mỗi lần mở màn. Sau khi Ngô Gia áp thay đổi, hồ sơ lịch sử không còn
cho biết trước đó là gì. Anh có cần lưu lại giá trị cũ thật không?
\end{ask}

% =============================================================================
\section{Huỷ yêu cầu cập nhật thông tin}
\rt{/portal/info-request/<int:req_id>/cancel}
\medskip

\mvao

Nút \emph{Huỷ} ở màn chi tiết.

\mai

Yêu cầu phải thuộc cửa hàng của tôi \textbf{và do chính tôi tạo}. Người khác cùng cửa hàng xem
được nhưng không huỷ được.

\mhien

Không hiện gì --- quay về màn chi tiết kèm băng thông báo.

\mloc

Không có ô nhập.

\mkiem{Huỷ}

\begin{kiembang}
1 & Yêu cầu thuộc cửa hàng của tôi và do tôi tạo & Đưa về danh sách \\
2 & Trạng thái đang là \emph{Nháp} hoặc \emph{Đã gửi} &
    Băng báo ``Chỉ có thể huỷ yêu cầu ở trạng thái Nháp/Đã gửi.'' \\
\end{kiembang}

\mghi

Chuyển trạng thái sang \textbf{Từ chối} và ghi lý do \emph{``User huỷ''}.

\mhoi

\begin{ask}
\textbf{Huỷ và bị từ chối trông giống nhau.} Hệ thống không có trạng thái \emph{Đã huỷ} nên
yêu cầu tự huỷ hiện đúng huy hiệu đỏ \emph{Từ chối} như khi Ngô Gia từ chối; chỉ khác dòng lý
do. Anh muốn tách thành trạng thái riêng không?
\end{ask}

% =============================================================================
\section{Lấy giá trị hiện tại của cửa hàng}
\rt{/portal/info-request/franchise/<int:fid>/values}
\medskip

\mvao

Không phải màn --- gọi ngầm từ form tạo yêu cầu mỗi lần người dùng đổi ô \emph{Loại thông tin},
để điền ô \emph{Giá trị hiện tại}.

\mai

Chỉ kiểm \textbf{một điều}: cửa hàng phải nằm trong danh sách tôi truy cập được.
\textbf{Không kiểm vai trò} --- khác với màn tạo yêu cầu vốn chỉ cho Chủ tiệm / Quản lý.

\mhien

Trả về đúng một giá trị: nội dung của trường tương ứng trên hồ sơ cửa hàng. Bốn loại thông tin
có sẵn ánh xạ (Địa chỉ, Số điện thoại, Email, Tên chủ cửa hàng); loại \emph{Khác} thì lấy theo
\textbf{tên trường người gọi truyền lên}.

\mloc

Hai tham số: loại thông tin và tên trường.

\mkiem{gọi}

\begin{kiembang}
1 & Cửa hàng nằm trong danh sách của tôi & Trả mã \rt{forbidden} \\
2 & Cửa hàng còn tồn tại & Trả mã \rt{not_found} \\
3 & Tên trường có thật trên hồ sơ cửa hàng & Trả chuỗi rỗng \\
\end{kiembang}

\mghi

\textbf{Không ghi gì.}

\mhoi

\begin{ask}
\textbf{Đường dẫn này đọc được BẤT KỲ trường nào của hồ sơ cửa hàng}, chỉ cần biết tên kỹ thuật
--- vì loại \emph{Khác} nhận thẳng tên trường do người gọi truyền lên và không có danh sách
trường cho phép. Nó lại \textbf{không kiểm vai trò}, nên một Nhân viên đọc được cả những số mà
màn Công nợ đã chặn họ (ví dụ số dư công nợ lưu trên hồ sơ cửa hàng). Dev sẽ đưa vào đợt chuẩn
hoá F1 --- ghi ở đây để BA nắm, \textbf{không cần lên task riêng}.
\end{ask}

% =============================================================================
\section{Màn Danh sách bài viết kiến thức}
\rt{/portal/knowledge}
\medskip

\mvao

Mục \emph{Kiến thức} trên thanh điều hướng. Tham số: \rt{page}, \rt{category_id},
\rt{tag_id}, \rt{keyword}, \rt{page_size}, \rt{notice}.

\mai

\begin{itemize}[leftmargin=*]
  \item Chỉ cần \textbf{đăng nhập}. Không cần chọn cửa hàng, \textbf{không} chặn vai trò,
        \textbf{không} lọc theo cửa hàng.
  \item Mọi tài khoản portal thấy \textbf{đúng cùng một tập bài viết} --- phân hệ duy nhất
        trong portal như vậy.
\end{itemize}

\mhien

\begin{hienbang}
Lưới bài viết &
  \rt{wujia.knowledge.article} &
  Bài \textbf{đang hiện trên portal} (đã phát hành, chưa lưu trữ, chưa hết hạn hiệu lực)
  \textbf{và} đã tới ngày phát hành (\rt{publish_date} rỗng hoặc đã qua) &
  Số thứ tự, rồi ngày phát hành mới nhất; \textbf{12} bài mỗi trang (12/24/48) &
  \uat{21}/27 bài đang hiện \\ \addlinespace
Trạng thái bài không hiện &
  \rt{state} &
  \uat{3} bài còn nháp · \uat{1} bài đã lưu trữ · \uat{2} bài \textbf{đã phát hành nhưng quá
  hạn hiệu lực} &
  --- &
  \uat{23} bài ở trạng thái đã phát hành \\ \addlinespace
Cột danh mục &
  \rt{wujia.knowledge.category} &
  Danh mục đang bật &
  Số thứ tự, rồi tên &
  \uat{9} danh mục --- trong đó có \textbf{tên trùng nhau} \\ \addlinespace
Khối thẻ &
  \rt{wujia.knowledge.tag} &
  Thẻ đang bật &
  Theo tên &
  \uat{6} thẻ \\ \addlinespace
Lượt xem &
  \rt{view_count} &
  Đếm mỗi lần mở màn chi tiết &
  --- &
  Bài cao nhất \uat{65} lượt \\
\end{hienbang}

\mloc

\begin{itemize}[leftmargin=*]
  \item \textbf{Từ khoá} --- khớp một phần theo \textbf{tiêu đề}, \textbf{tóm tắt} hoặc
        \textbf{toàn văn nội dung}. Đây là ô tìm sâu nhất trong portal.
  \item \textbf{Danh mục} và \textbf{Thẻ} --- chọn một cái mỗi loại. Nếu danh mục hay thẻ vừa bị
        Ngô Gia lưu trữ (hoặc số gửi lên không đọc được) thì màn \textbf{quay về danh sách gốc}
        kèm băng báo ``danh mục/thẻ không còn'' --- cố ý không im lặng bỏ lọc.
  \item \textbf{Số bài mỗi trang} --- \textbf{12, 24, 48} (bội của 3 cho lưới 3 cột), khác bậc
        10/20/50 của các màn khác.
\end{itemize}

\mkiem{một bài viết}

Không kiểm gì ở màn danh sách.

\mghi

\textbf{Không ghi gì.}

\mhoi

\begin{ask}
\begin{enumerate}[leftmargin=*]
  \item \textbf{Kiến thức không phân biệt cửa hàng và vai trò.} Mọi nhân viên của mọi cửa hàng
        đọc được mọi bài. Nếu sau này có tài liệu chỉ dành cho Chủ tiệm thì phải làm thêm.
  \item \textbf{Danh mục đang trùng tên trên UAT} --- \emph{Vận hành} và \emph{Sản phẩm} mỗi cái
        xuất hiện hai lần (hai bản ghi khác nhau). Thanh lọc bên trái vì thế hiện hai dòng
        giống hệt. Cần dọn dữ liệu.
  \item \textbf{Bài hết hạn không biến mất ngay.} Cờ ``đang hiện'' được lưu sẵn và chỉ tính lại
        bằng tác vụ nền chạy hằng ngày, nên bài đặt hạn 10 giờ sáng vẫn đọc được tới lần chạy
        kế tiếp. Anh chấp nhận độ trễ này chứ?
  \item \textbf{Lượt xem đếm mỗi lần mở}, không chống bấm lặp --- một người bấm lại 10 lần thì
        cộng 10. Con số \uat{65} lượt trên UAT nên hiểu theo nghĩa đó.
\end{enumerate}
\end{ask}

% =============================================================================
\section{Màn Chi tiết bài viết}
\rt{/portal/knowledge/<string:slug>}
\medskip

\mvao

Bấm một thẻ bài ở màn danh sách, hoặc từ khối \emph{Kiến thức} trên Trang chủ, hoặc từ ô gợi ý
tìm kiếm. Đường dẫn dùng \textbf{tên rút gọn} của bài chứ không dùng số --- ví dụ
\rt{/portal/knowledge/checklist-mo-cua-hang-buoi-sang}.

\mai

Bài phải \textbf{đang hiện trên portal} và đã tới ngày phát hành. Không thoả --- kể cả khi tên
rút gọn gõ sai --- thì về danh sách kèm băng ``bài viết không còn''.

\mhien

\begin{hienbang}
Nội dung bài &
  \rt{wujia.knowledge.article} &
  Tiêu đề, tóm tắt, ngày phát hành, danh mục, thẻ, toàn văn nội dung &
  --- & --- \\ \addlinespace
Tệp đính kèm &
  \rt{ir.attachment} &
  Tệp gắn vào bài, tải qua đường dẫn có kiểm quyền &
  --- & --- \\ \addlinespace
Bài liên quan &
  \rt{wujia.knowledge.article} &
  Bài \textbf{cùng danh mục}, đang hiện, trừ chính bài đang đọc &
  Ngày phát hành mới nhất; \textbf{tối đa 4} bài &
  --- \\
\end{hienbang}

\mloc

Không có ô lọc.

\mkiem{một bài liên quan}

Không kiểm gì --- mở sang màn chi tiết của bài đó, và màn đó kiểm lại từ đầu.

\mghi

\textbf{Có ghi} --- đây là màn chỉ-đọc duy nhất trong portal có ghi dữ liệu: mỗi lần mở,
\textbf{lượt xem của bài cộng thêm 1}. Phép cộng chạy thẳng dưới cơ sở dữ liệu nên hai người mở
cùng lúc không ghi đè nhau. Không ghi ai đã đọc bài --- \textbf{không có} khái niệm ``đã đọc''
cho Kiến thức.

\mhoi

\begin{ask}
\begin{enumerate}[leftmargin=*]
  \item \textbf{Không biết ai đã đọc bài nào.} Nếu Ngô Gia cần chắc rằng cửa hàng đã đọc tài
        liệu bắt buộc thì phải làm thêm phần xác nhận đã đọc.
  \item \textbf{Đường dẫn dùng tên rút gọn.} Đổi tiêu đề bài \textbf{không} đổi tên rút gọn, nên
        liên kết cũ vẫn sống --- nhưng tên rút gọn sẽ không còn khớp tiêu đề.
\end{enumerate}
\end{ask}

% =============================================================================
\section{Tải tệp đính kèm của bài viết}
\rt{/portal/knowledge/<string:slug>/attachment/<int:att_id>}
\medskip

\mvao

Bấm tên tệp ở màn chi tiết bài.

\mai

\begin{kiembang}
1 & Bài đang hiện trên portal & Về danh sách kèm băng ``bài viết không còn'' \\
2 & Tệp thuộc đúng bài đó & Từ chối truy cập \\
\end{kiembang}

Ở đây \textbf{cố ý không} kiểm theo cửa hàng --- kiến thức là tài liệu chung cho mọi tài khoản
portal.

\mhien

Trả về tệp ở chế độ \textbf{tải xuống}.

\mloc

Không có.

\mkiem{tên tệp}

Xem bảng trên.

\mghi

\textbf{Không ghi gì} --- tải tệp không cộng lượt xem.

\mhoi

\begin{ask}
Không có điểm cần chốt.
\end{ask}

% =============================================================================
\section{Gợi ý tìm kiếm kiến thức}
\rt{/portal/knowledge/search}
\medskip

\mvao

Không phải màn --- gọi ngầm khi người dùng đang gõ trong ô tìm kiếm, sau khi ngừng gõ khoảng
0,3 giây.

\mai

Chỉ cần đăng nhập. Cùng phạm vi với màn danh sách: mọi bài đang hiện.

\mhien

Trả về danh sách gợi ý, mỗi dòng gồm tiêu đề, danh mục và đường dẫn tới bài. Tìm theo tiêu đề,
tóm tắt và toàn văn nội dung.

\mloc

\begin{itemize}[leftmargin=*]
  \item \textbf{Từ khoá} --- \textbf{dưới 2 ký tự thì không gọi}, trả về rỗng.
  \item \textbf{Số gợi ý} --- mặc định 10, tối đa \textbf{20}; số lạ tự về 10.
\end{itemize}

\mkiem{gõ phím}

Không kiểm gì ngoài độ dài từ khoá.

\mghi

\textbf{Không ghi gì} --- gợi ý không cộng lượt xem.

\mhoi

\begin{ask}
\textbf{Chưa có giới hạn số lần gọi.} Gõ nhanh và liên tục sẽ tạo nhiều truy vấn tìm toàn văn;
với 1.500 tài khoản cần theo dõi khi chạy thật.
\end{ask}

% =============================================================================
\section{Màn Danh sách thông báo}
\rt{/portal/notification} · \rt{/portal/notification/results}
\medskip

\mvao

Mục \emph{Thông báo} trên thanh điều hướng, hoặc dòng \emph{Xem tất cả} trong hộp chuông.
Đường dẫn \rt{/results} không phải màn riêng --- là phần kết quả gọi ngầm khi đổi bộ lọc.
Tham số: \rt{page}, \rt{limit}, \rt{type_id}, \rt{keyword}, \rt{priority},
\rt{read_status} (hoặc \rt{unread} kiểu cũ), \rt{date_from}, \rt{date_to}.

\mai

\begin{itemize}[leftmargin=*]
  \item Phải đăng nhập. \textbf{Không chặn vai trò.}
  \item Thấy thông báo \textbf{gửi toàn hệ} \emph{hoặc} \textbf{gửi đích danh một trong các cửa
        hàng của tôi}. Chưa chọn cửa hàng thì vẫn thấy phần gửi toàn hệ.
  \item Danh sách là \textbf{lịch sử}: gồm cả thông báo \textbf{đã hết hiệu lực} (hộp chuông thì
        không).
  \item Thực tế UAT: \uat{3} thông báo, \uat{2} đang phát hành --- cả hai đều gửi toàn hệ;
        \uat{0} thông báo gửi đích danh cửa hàng đang hiện ra portal. Nghĩa là luật lọc theo
        cửa hàng \textbf{chưa được dữ liệu UAT kiểm chứng}.
\end{itemize}

\mhien

\begin{hienbang}
Danh sách thông báo &
  \rt{wujia.notification} &
  Đã phát hành ra portal (\rt{is_published_portal = true}), \textbf{đã tới giờ phát hành},
  và đúng đối tượng nhận &
  \textbf{Ghim lên đầu}, rồi ngày phát hành mới nhất; \textbf{10} thông báo mỗi trang
  (10/20/50) &
  \uat{2} đang phát hành · \uat{1} chưa phát hành · \uat{0} được ghim \\ \addlinespace
Huy hiệu loại &
  \rt{wujia.notification.type} &
  Năm loại kèm màu riêng: Khẩn cấp · Thông báo chung · Khuyến mãi · Hệ thống · Khác &
  Theo số thứ tự &
  \uat{5} loại đang bật \\ \addlinespace
Huy hiệu mức độ &
  \rt{priority} &
  Ba nhãn: \emph{Thông thường} · \emph{Quan trọng} · \emph{Cần làm} &
  --- & --- \\ \addlinespace
Dấu đã đọc &
  \rt{wujia.notification.read} &
  Đánh dấu theo bộ ba \textbf{người + thông báo + cửa hàng đang chọn}. Chưa chọn cửa hàng thì
  bỏ điều kiện cửa hàng khi tra. &
  --- &
  \uat{6} dấu cũ chưa gắn cửa hàng \\ \addlinespace
Huy hiệu hết hiệu lực &
  \rt{expired_date} &
  Thông báo quá hạn vẫn nằm trong danh sách nhưng được đánh dấu riêng &
  --- &
  \uat{0} thông báo đã hết hiệu lực \\
\end{hienbang}

\mloc

\begin{itemize}[leftmargin=*]
  \item \textbf{Từ khoá} --- tiêu đề, \textbf{số công văn}, hoặc tóm tắt.
  \item \textbf{Loại} và \textbf{Mức độ} --- mỗi thứ chọn một; giá trị lạ bị bỏ qua.
  \item \textbf{Trạng thái đọc} --- \emph{Tất cả} / \emph{Đã đọc} / \emph{Chưa đọc}. Chọn
        \emph{Chưa đọc} thì hệ thống \textbf{tự loại luôn thông báo đã hết hiệu lực} --- quá hạn
        không bị tính là còn nợ đọc.
  \item \textbf{Khoảng ngày} --- lọc theo \emph{ngày phát hành}, quy đổi đúng múi giờ tài khoản.
        Ngày ngược thì \textbf{chặn truy vấn}, giữ chữ đã gõ và báo tại thanh lọc --- cố ý không
        bỏ lọc rồi trả về toàn bộ.
  \item \textbf{Số dòng mỗi trang} --- 10, 20, 50.
\end{itemize}

\mkiem{một thông báo}

Không kiểm gì ở danh sách.

\mghi

\textbf{Không ghi gì} --- xem danh sách không đánh dấu đã đọc.

\mhoi

\begin{ask}
\begin{enumerate}[leftmargin=*]
  \item \textbf{``Đã đọc'' tính theo từng cửa hàng.} Người phụ trách 3 cửa hàng phải đọc cùng
        một thông báo toàn hệ \textbf{ba lần} mới sạch chuông. Anh muốn giữ vậy, hay đọc một
        lần là xong cho mọi cửa hàng?
  \item \uat{6} dấu đã đọc cũ \textbf{chưa gắn cửa hàng} nên không khớp cửa hàng nào --- chúng
        chỉ có tác dụng khi người dùng \textbf{chưa chọn} cửa hàng. Cần dọn dữ liệu.
  \item \textbf{Thẻ đếm chưa đọc trên Trang chủ tính khác màn này} --- Trang chủ bỏ qua điều
        kiện hết hiệu lực và điều kiện đã tới giờ phát hành, lại trừ \textbf{mọi} dấu đã đọc.
        Đây là chênh lệch đã nêu ở chương 4; nhắc lại để BA gom vào một quyết định.
  \item \textbf{UAT chưa có thông báo gửi đích danh cửa hàng nào đang phát hành} --- cần dựng ca
        mẫu trước khi nghiệm thu phần lọc theo cửa hàng.
\end{enumerate}
\end{ask}

% =============================================================================
\section{Màn Chi tiết thông báo}
\rt{/portal/notification/<int:notification_id>}
\medskip

\mvao

Bấm một dòng ở danh sách hoặc một dòng trong hộp chuông.

\mai

Thông báo phải \textbf{đã phát hành, đã tới giờ, đúng đối tượng nhận}. Thông báo
\textbf{đã hết hiệu lực vẫn mở được} --- cố ý, để tra lại lịch sử. Không thoả thì về danh sách.

\mhien

\begin{hienbang}
Nội dung &
  \rt{wujia.notification} &
  Tiêu đề, số công văn, loại, mức độ, ngày phát hành, toàn văn nội dung &
  --- & --- \\ \addlinespace
Tệp đính kèm &
  \rt{ir.attachment} &
  Tệp gắn vào thông báo, tải qua đường dẫn có kiểm quyền &
  --- &
  \uat{0}/3 thông báo có tệp \\ \addlinespace
Ghi chú nội bộ &
  --- &
  \textbf{Không hiện} --- ô ghi chú của Ngô Gia không ra portal &
  --- & --- \\
\end{hienbang}

\mloc

Không có ô lọc.

\mkiem{mở màn}

\begin{kiembang}
1 & Thông báo đã phát hành, đã tới giờ, đúng đối tượng & Về danh sách \\
2 & Đã chọn được cửa hàng &
    \textbf{Vẫn cho đọc nội dung}, chỉ \textbf{không ghi} dấu đã đọc --- tránh tạo dấu không
    gắn cửa hàng nào \\
\end{kiembang}

\mghi

\begin{itemize}[leftmargin=*]
  \item Lần đầu mở tại một cửa hàng: tạo một dấu đã đọc cho bộ ba \textbf{người + thông báo +
        cửa hàng}, ghi cả \emph{lần đọc đầu} và \emph{lần mở gần nhất}.
  \item Mở lại: \textbf{chỉ cập nhật lần mở gần nhất}, giữ nguyên lần đọc đầu.
  \item Chưa chọn cửa hàng: \textbf{không ghi gì}, và chuông vẫn coi là chưa đọc.
\end{itemize}

\mhoi

\begin{ask}
\textbf{Mở thông báo khi chưa chọn cửa hàng thì không được tính là đã đọc} --- người dùng đọc
xong vẫn thấy chuông báo đỏ. Cần nói rõ trên màn hay tự chọn cửa hàng giúp họ?
\end{ask}

% =============================================================================
\section{Hộp chuông trên đầu trang}
\rt{/portal/notification/recent} · \rt{/portal/notification/unread-count}
\medskip

\mvao

Không phải màn --- hai đường dẫn phục vụ biểu tượng chuông có mặt trên \textbf{mọi trang}
portal: một cái nạp nội dung khi bấm mở, một cái chỉ trả về con số trên chuông.

\mai

Phải đăng nhập. Phạm vi giống màn danh sách, nhưng \textbf{chỉ thông báo còn hiệu lực} --- quá
hạn thì rụng khỏi chuông dù vẫn nằm trong lịch sử.

\mhien

\begin{hienbang}
Danh sách trong hộp &
  \rt{wujia.notification} &
  Thông báo \textbf{còn hiệu lực} đúng đối tượng nhận &
  Ghim lên đầu, rồi mới nhất; \textbf{đúng 5} dòng &
  \uat{2} thông báo còn hiệu lực trên UAT \\ \addlinespace
Số trên chuông &
  \rt{wujia.notification.read} &
  \textbf{Số thông báo còn hiệu lực trừ đi số dấu đã đọc} của tôi tại cửa hàng đang chọn; sàn ở
  0 &
  --- &
  Toàn hệ có \uat{18} dấu đã đọc, trong đó \uat{6} dấu chưa gắn cửa hàng \\
\end{hienbang}

Hai đường dẫn này chạy trên mọi trang nên được viết để \textbf{chỉ đếm, không kéo dữ liệu về} ---
điều kiện ``còn hiệu lực'' được đẩy xuống thành câu hỏi phụ trong cơ sở dữ liệu thay vì nạp
danh sách mã ra rồi so ở tầng ứng dụng.

\mloc

Không có ô nhập.

\mkiem{mở chuông}

Không kiểm gì ngoài đăng nhập.

\mghi

\textbf{Không ghi gì} --- mở chuông không đánh dấu đã đọc.

\mhoi

\begin{ask}
\textbf{Hộp chuông cố định 5 dòng} và không có nút ``xem thêm'' trong hộp --- phải sang màn danh
sách. Anh xác nhận.
\end{ask}

% =============================================================================
\section{Đánh dấu đã đọc}
\rt{/portal/notification/mark-all-read} · \rt{/portal/notification/mark-read}
\medskip

\mvao

Nút \emph{Đánh dấu tất cả đã đọc} trong hộp chuông (đường dẫn thứ nhất); đường dẫn thứ hai giữ
lại cho các màn cũ gửi lên một danh sách mã cụ thể.

\mai

Phải đăng nhập \textbf{và phải đang chọn một cửa hàng}. Chưa chọn thì cả hai trả về lỗi
\emph{``Vui lòng chọn cửa hàng trước khi thao tác''} và \textbf{không ghi gì}.

\mhien

Không hiện gì --- trả về số dấu vừa tạo và số chưa đọc còn lại.

\mloc

Đường dẫn thứ hai nhận một danh sách mã thông báo; mã không đọc được thì trả lỗi
\rt{invalid_ids}.

\mkiem{Đánh dấu tất cả đã đọc}

\begin{kiembang}
1 & Đang chọn một cửa hàng & ``Vui lòng chọn cửa hàng trước khi thao tác.'' \\
2 & (đường dẫn thứ hai) Mã gửi lên đều là số & \rt{invalid_ids} \\
3 & (đường dẫn thứ hai) Thông báo \textbf{đúng đối tượng nhận} &
    Mã lạ bị \textbf{lặng lẽ bỏ qua} --- không báo lỗi, không đánh dấu \\
\end{kiembang}

\mghi

\begin{itemize}[leftmargin=*]
  \item \textbf{Đánh dấu tất cả}: tạo dấu đã đọc cho \textbf{mọi thông báo còn hiệu lực} chưa
        có dấu, tại \textbf{cửa hàng đang chọn}. \textbf{Không} ghi mốc ``lần mở gần nhất'' ---
        người dùng chưa thực sự mở nội dung.
  \item \textbf{Đánh dấu theo danh sách}: cùng cách, nhưng có ghi cả mốc lần mở.
  \item Cả hai đều \textbf{chạy lại được nhiều lần không sinh trùng} --- bộ ba người + thông báo
        + cửa hàng là duy nhất.
  \item \textbf{Chỉ đánh dấu cho cửa hàng đang chọn}: người ba cửa hàng bấm \emph{tất cả} vẫn
        còn chuông đỏ ở hai cửa hàng kia.
\end{itemize}

\mhoi

\begin{ask}
\textbf{``Đánh dấu tất cả'' không phải là tất cả} --- chỉ là tất cả \emph{tại cửa hàng đang
chọn} và chỉ những thông báo \emph{còn hiệu lực}. Nhãn nút nên nói rõ điều đó.
\end{ask}

% =============================================================================
\section{Tải tệp đính kèm của thông báo}
\rt{/portal/notification/<int:notification_id>/attachment/<int:attachment_id>}
\medskip

\mvao

Bấm tên tệp ở màn chi tiết thông báo.

\mai

\begin{kiembang}
1 & Thông báo đã phát hành, đã tới giờ, \textbf{đúng đối tượng nhận của tôi} & Báo không tìm thấy \\
2 & Tệp thuộc đúng thông báo đó &
    Từ chối truy cập --- đây chính là lý do portal không dùng đường dẫn tải tệp mặc định của
    Odoo, vì đường dẫn đó chỉ cần biết số tệp là tải được \\
\end{kiembang}

\mhien

Trả về tệp ở chế độ \textbf{tải xuống}.

\mloc

Không có.

\mkiem{tên tệp}

Xem bảng trên.

\mghi

\textbf{Không ghi gì.}

\mhoi

\begin{ask}
Không có điểm cần chốt.
\end{ask}
